\documentclass[10pt]{article}
\usepackage[letterpaper,margin=1in]{geometry}
\usepackage{mathptmx}
\usepackage[utf8]{inputenc}
\usepackage{url}
\pagestyle{empty}
\setlength{\parindent}{0pt}
\setlength{\parskip}{6pt}
\begin{document}

\noindent
Kailin Yan\\
Department of Computer Science\\
China University of Petroleum-Beijing at Karamay\\
Karamay, China\\
\texttt{2024015943@st.cupk.edu.cn}\\[10pt]
\noindent 8 October 2026\\[10pt]

\noindent Dear Editorial Office,\\[6pt]

\noindent Thank you for the technical assessment of our submission
(unsubmitted 08-Oct-2026). We have addressed each item as follows and
resubmit.

\medskip
\noindent\textbf{Manuscript:} ``From Embodied Experience to Persistent
Experience-Derived State: An Auditable Unified-Graph Architecture and a
Preregistered Boundary Analysis of Activation-Based Access''

\medskip
\noindent\textbf{1.\ Institutional email and profile links (single
author).} Author: Kailin Yan. Institutional email:
\texttt{2024015943@st.cupk.edu.cn} (China University of
Petroleum-Beijing at Karamay --- student address on the university's
st.cupk.edu.cn domain; identical to the address on the manuscript title
page and the submission portal account). Profile link: ORCID
0009-0009-6762-2684 (\url{https://orcid.org/0009-0009-6762-2684}). The
author is a degree-seeking student without a personal institutional web
page; ORCID is provided as the persistent identifier.

\medskip
\noindent\textbf{2.\ Relevance to the journal.} The journal's scope
explicitly includes cognitive modeling, the design and experimental
evaluation of agent architectures, and embodied artificial intelligence.
This paper contributes on exactly those axes: (i) an auditable
unified-graph cognitive architecture (FAS) in which episodic, semantic
and procedural knowledge share one structure with a single
evidence-weighted write-back pathway; (ii) preregistered, fully archived
behavioral experiments in a deterministic embodied sandbox --- including
a preregistered rejection of the architecture's own headline hypothesis
(spreading activation showed no incremental value over flat retrieval);
and (iii) a boundary analysis, with parameter sweeps, of when
activation-based access does and does not pay. JETAI has historically
published the cognitive-architecture line this work builds on (Soar,
ACT-R-inspired designs) and current embodied-agent research; the
methodology (preregistration plus reproducibility package) is offered as
a contribution to the experimental standards of that community.

\medskip
\noindent\textbf{3.\ Technical-check items.}
\begin{itemize}
\item Table 3 is now referenced in the text (Section 5.1: ``Table 3
  reports the full grid.'').
\item The title page now includes the keyword list.
\item The affiliation in the manuscript files now matches the submission
  metadata exactly: ``Department of Computer Science, China University of
  Petroleum-Beijing at Karamay, Karamay, China''.
\item The manuscript body has been reduced to 4,982 words (excluding
  tables, figures, captions, algorithms and references). The relocated
  passages were moved verbatim --- no wording or number altered --- into
  the supplementary file, new Section S19 (subsections S19.1--S19.13),
  where each carries a ``moved from Section X'' heading; the main text
  retains summaries with pointers to S19.
\end{itemize}

\medskip
\noindent All data and code are deposited (Zenodo, DOI
10.5281/zenodo.23196280); the anonymized manuscript withholds the DOI
per double-blind policy.

\bigskip
\noindent Sincerely,\\[14pt]
\noindent Kailin Yan\\
\noindent Department of Computer Science\\
\noindent China University of Petroleum-Beijing at Karamay, Karamay, China\\
\noindent \texttt{2024015943@st.cupk.edu.cn}

\end{document}
